% Part: sets-functions-relations
% Chapter: sets
% Section: russells-paradox

\documentclass[../../../include/open-logic-section]{subfiles}


\begin{document}

\olfileid[tr]{sfr}{set}{rus}
\olsection{Russell Paradoksu}

Kaplamsallık, $\phi(x)$ özelliğini taşıyan $x$ değerlerinin \emph{kümesi} için
$\Setabs{x}{\phi(x)}$ gösterimini kullanmamıza olanak verir. Ancak
kaplamsallığın \emph{gerçekte} izin verdiği düşünce yalnızca şudur:
elemanları tam olarak $\phi$ özelliğini taşıyan nesneler olan bir küme
\emph{varsa}, bu tanımı karşılayan yalnızca bir küme vardır. Başka bir deyişle,
belirli bir $\phi$ seçildikten sonra $\Setabs{x}{\phi(x)}$ kümesi \emph{varsa}
tektir.

Ne var ki bu koşul önemlidir!{} En önemlisi, her özellik \emph{küme
oluşturmaya} elverişli değildir. Yani bazı özellikler küme tanımlamaz. Hepsi
tanımlasaydı doğrudan çelişkilere varırdık. Bunun en ünlü örneği Russell
paradoksudur.

Kümeler başka kümelerin elemanları olabilir---örneğin bir $A$ kümesinin
kuvvet kümesi, kümelerden oluşur. Dolayısıyla bir kümenin başka bir kümenin
elemanı olup olmadığını sormak veya araştırmak anlamlıdır. Bir küme kendisinin
elemanı olabilir mi? Küme düşüncesindeki hiçbir şey bunu
dışlıyor gibi görünmez. Örneğin \emph{bütün} kümeler bir nesneler topluluğu
oluşturuyorsa bunların tek bir kümede---bütün kümelerin kümesinde---
toplanabileceği düşünülebilir. Bütün kümelerin kümesi de kendisi bir küme
olduğuna göre kendi elemanı olurdu.

Russell paradoksu, bir kümenin \emph{kendi kendisinin elemanı olmaması}
özelliğini ele aldığımızda ortaya çıkar. Kendi kendisinin elemanı olmayan bütün kümelerden
oluşan bir küme
bulunduğunu varsayarsak ne olur?
\[
R = \Setabs{x}{x \notin x}
\]
Böyle bir $R$ var mıdır? Böyle bir kümenin var olmadığını
ispatlayabileceğimiz ortaya çıkar.

\begin{thm}[Russell Paradoksu]\ollabel{thm:russells-paradox}
	$R = \Setabs{x}{x \notin x}$ eşitliğini sağlayan hiçbir $R$ kümesi yoktur.
\end{thm}

\begin{proof}
$R = \Setabs{x}{x \notin x}$ eşitliğini sağlayan bir $R$ kümesi varsa,
$R \in R$ ancak ve ancak $R \notin R$ olur; bu da bir çelişkidir.
\end{proof}

\begin{tagblock}{novice}
\begin{explain}
Bu ispatı daha yavaş inceleyelim. $R$ varsa $R \in R$ olup olmadığını sormak
anlamlıdır. Önce gerçekten $R \in R$ olduğunu varsayalım. $R$, kendi kendisini
eleman olarak içermeyen bütün kümelerin kümesi diye tanımlanmıştı. Dolayısıyla
$R \in R$ ise $R$'nin kendisi, $R$'yi tanımlayan özelliği taşımaz. Oysa
yalnızca bu özelliği taşıyan kümeler $R$'de bulunur; öyleyse $R$, $R$'nin
elemanı olamaz, yani $R \notin R$ olur. Fakat $R$ aynı anda hem kendisinin elemanı hem de kendisinin
elemanı olmayan bir küme olamaz; dolayısıyla bir çelişkiye ulaştık.

$R \in R$ varsayımı çelişkiye götürdüğüne göre $R \notin R$ sonucunu elde
ederiz. Ancak bu da bir çelişkiye götürür!{} Çünkü $R \notin R$ ise $R$'nin
kendisi, $R$'yi tanımlayan özelliği taşır; o hâlde kendi kendisinin elemanı
olmayan diğer tüm kümeler gibi $R$ de kendisinin elemanı olurdu. Yine $R$
aynı anda hem kendisinin elemanı olmayan hem de kendisinin elemanı olan bir
küme olamaz.
\end{explain}
\end{tagblock}

\begin{digress}
Russell paradoksuna düşmeyen, yani
$R = \Setabs{x}{x \notin x}$ kümesinin var olduğu yönündeki \emph{tutarsız}
savı ileri sürmekten kaçınan bir küme kuramını nasıl kurabiliriz? Kümelerin ne
zaman var olduğunu (ve ne zaman var olmadığını) söylemek için çok kesin
koşullar veren aksiyomlar koymamız gerekir.
	
Bu bölümde ana hatları çizilen küme kuramı bunu yapmaz. Kuram
\emph{gerçekten naiftir}. Yalnızca kümelerin kaplamsallığa uyduğunu ve bazı
kümeler varsa bunların birleşimini, kesişimini vb. oluşturabileceğinizi
söyler. Küme kuramını bundan daha titiz biçimde geliştirmek mümkündür.
\oliflabeldef{cumul:::part}{Bu titizlik, \olref[cumul][][]{part}.~Kısım için
saklanacaktır. Şimdilik naif biçimde ilerleyecek ve çelişkilerden dikkatle
kaçınmaya çalışacağız.}{}
\end{digress}


\end{document}
